% =====================================================================
%  9. CONCLUSIONES
%  Responsable: INTEGRANTE 5
%  Extensión objetivo: 1,5 páginas — aprox. 550 palabras
%
%  DOS REGLAS QUE NO SE NEGOCIAN:
%   1. UNA CONCLUSIÓN POR CADA OBJETIVO ESPECÍFICO de la sección 2, en el
%      mismo orden. Si hay cinco objetivos, hay cinco conclusiones.
%   2. NO se introduce información nueva ni citas nuevas. Todo lo que se
%      concluye ya debe estar demostrado en las secciones 4 a 8.
% =====================================================================

\chapter{Conclusiones}

% TODO (conclusión 1 — responde al objetivo específico 1: describir la
%   unidad minera y su entorno):
%   La Mina Marlin fue una operación mixta de oro y plata sobre un depósito
%   epitermal de baja sulfuración, implantada en territorio de los pueblos
%   maya-mam y maya-sipakapense. Cierra señalando que la condición indígena
%   del territorio es el factor que determina todo el perfil de riesgo social.

% TODO (conclusión 2 — responde al objetivo 2: alcance, criterios y
%   metodología de la auditoría):
%   La HRIA fue una auditoría de desempeño en derechos humanos, no de
%   sistema de gestión, ejecutada por una firma independiente sobre criterios
%   internacionales y con alcance temporal y espacial declarados, incluidas
%   sus exclusiones.

% TODO (conclusión 3 — responde al objetivo 3: clasificar los hallazgos):
%   Indica cuántos hallazgos se clasificaron y en qué categorías, y cuál es
%   el eje más comprometido. Nombra el hallazgo de mayor gravedad (ausencia
%   de consulta previa) y por qué lo es.

% TODO (conclusión 4 — responde al objetivo 4: evaluar la independencia y
%   las limitaciones):
%   Toma posición argumentada sobre la tensión entre el financiamiento por
%   la empresa auditada y el diseño institucional que buscó mitigarlo, y
%   sobre el efecto del sesgo de participación en la validez de los hallazgos.

% TODO (conclusión 5 — responde al objetivo 5: proponer el plan PHVA):
%   Sintetiza qué aporta el plan propuesto y sobre qué estándares se apoya.

% TODO (conclusión de cierre, opcional pero recomendable):
%   Una reflexión final sobre la utilidad y los límites de la auditoría
%   socioambiental voluntaria como herramienta de control, y su relación con
%   la fiscalización estatal. Sin retórica: una idea, bien dicha.

\begin{enumerate}
  \item \pendiente{conclusión 1 — sobre la unidad minera}
  \item \pendiente{conclusión 2 — sobre el alcance y la metodología de la auditoría}
  \item \pendiente{conclusión 3 — sobre la clasificación de hallazgos}
  \item \pendiente{conclusión 4 — sobre la independencia y las limitaciones}
  \item \pendiente{conclusión 5 — sobre el plan de mejora continua}
\end{enumerate}
